%! Function Operations — Unit 1, Lesson 1.2 — Homework (ANSWER KEY)
% THE LESSON'S GRADED PRACTICE and the last component in the packet — exactly TWO pages, started
% in the Homework Launch (the last 3 minutes of the period) and finished at home. Every context
% here is new: the notes used a bakery and a food truck; the AP Practice used a school concert
% and a bike shop; the homework uses a print shop and a lawn-care service.
%   Page 1 — Part A (at a point, 1–2) and Part B (building the formula and its domain, 3–5).
%   Page 2 — Part C (in context, 6–7) and Part D (the error and the headline, 8–9).
% The diagnostic item is problem 4(b): the quotient that cancels, met alone. Problem 9 asks for
% the lesson's headline in writing.
%
% PAGE LOCKSTEP with the other half of the pair: work blocks are byte-identical, prose answers
% are \writespace{H}{…} (fixed height both ways), and each \blank sits at the END of a line so
% its \ans cannot rewrap anything. The single explicit \newpage fixes the break in both files.
\documentclass[10pt]{article}
\usepackage{precalculus-article}
\usepackage{precalculus-key}
\newcommand{\TallMath}[1]{$\displaystyle #1\rule[-1.4em]{0pt}{3.2em}$}

\begin{document}

\pageheader{Unit 1, Lesson 1.2}{Homework --- Answer Key}

\vspace{0.1in}

\boxguard[8]
\begin{remindbox}
\small
\textbf{This is your graded homework.} Start it in class; finish it on your own by the due date
on your cover sheet. Show your work: put each function's whole formula in parentheses before you
subtract, and settle the \emph{domain} from the original functions before you simplify anything.
\end{remindbox}

\vspace{0.1in}

\begin{headlinebox}{lilac}
\textbf{\color{plum}Part A --- Combining at a point}
\end{headlinebox}

\begin{enumerate}[leftmargin=1.9em, itemsep=8pt]

\item \begin{minipage}[t]{0.48\linewidth}
      \vspace{0pt}
      Use the table of values for $r$ and $s$.

      \smallskip
      $(r+s)(2) = $ \ans{$7$}

      \smallskip
      $(rs)(1) = $ \ans{$15$}

      \smallskip
      $(r-s)(0) = $ \ans{$3$}
      \end{minipage}\hfill
      \begin{minipage}[t]{0.46\linewidth}
      \vspace{0pt}
      \centering
      {\renewcommand{\arraystretch}{1.3}
      \begin{tabular}{|c|c|c|c|c|c|}
      \hline
      $x$    & $-1$ & 0 & 1 & 2 & 3 \\ \hline
      $r(x)$ & 6    & 4 & 3 & $-2$ & $-5$ \\ \hline
      $s(x)$ & $-3$ & 1 & 5 & 9 & 13 \\ \hline
      \end{tabular}}
      \end{minipage}

\item From the same table: \quad $(r-s)(3) = $ \ans{$-18$} \qquad $(s-r)(3) = $ \ans{$18$}

      How are the two answers related, and why does that always happen?
      \writespace{0.9cm}{They are opposites: $-18$ and $18$. Subtracting in the other order negates every output, so $(s-r)(x) = -(r-s)(x)$ always.}

\end{enumerate}

\vspace{0.1in}

\begin{headlinebox}{lilac}
\textbf{\color{plum}Part B --- Building the formula, and its domain}
\end{headlinebox}

\begin{enumerate}[leftmargin=1.9em, itemsep=8pt, start=3]

\item Let $f(x) = 3x - 4$ and $g(x) = x + 9$.

      $(f+g)(x) = $ \ans{$4x + 5$} \qquad $(f-g)(x) = $ \ans{$2x - 13$}

      $(g-f)(x) = $ \ans{$-2x + 13$} \qquad $(fg)(x) = $ \ans{$3x^2 + 23x - 36$}

\item Let $f(x) = x^2 - 16$ and $g(x) = x - 4$.

      \textbf{(a)} Write $\left(\dfrac{f}{g}\right)(x)$ and simplify it.

\begin{work}
  \frac{x^2 - 16}{x - 4} &= \frac{(x-4)(x+4)}{x-4} \\
                         &= x + 4
\end{work}

      Simplified form: \ans{$x + 4$}

      \textbf{(b)} State the domain of $\dfrac{f}{g}$, and explain why it is not ``all real
      numbers.''
      \writespace{1.2cm}{Every real number except $4$. The quotient was never defined at $x = 4$, because $g(4) = 0$; cancelling the $(x-4)$ afterwards cannot put that input back.}

\item Let $f(x) = \dfrac{1}{x + 3}$ and $g(x) = 2x$.

      Domain of $f + g$: \ans{every real number except $-3$}

      Domain of $\dfrac{f}{g}$: \ans{every real number except $-3$ and $0$}

\end{enumerate}

\newpage

\begin{headlinebox}{lilac}
\textbf{\color{plum}Part C --- Combinations in context}
\end{headlinebox}

\begin{enumerate}[leftmargin=1.9em, itemsep=8pt, start=6]

\item A print shop fills an order of $x$ posters. Its ink cost is $I(x) = 0.75x$ dollars and its
      paper-and-setup cost is $P(x) = 0.50x + 40$ dollars.

      \textbf{(a)} Write the total cost $T(x) = (I + P)(x)$ in simplified form.

\begin{work}
  T(x) &= 0.75x + (0.50x + 40) \\
       &= 1.25x + 40
\end{work}

      $T(x) = $ \ans{$1.25x + 40$}

      \textbf{(b)} Find $T(200)$ and say what it means, with units. \quad $T(200) = $ \ans{$290$}

\item A lawn-care service mows $x$ lawns in a day and can handle at most 12. Its revenue is
      $R(x) = 35x$ dollars and its cost is $C(x) = 12x + 80$ dollars.

      \textbf{(a)} Write the profit $P(x) = (R - C)(x)$ in simplified form.

\begin{work}
  P(x) &= 35x - (12x + 80) \\
       &= 23x - 80
\end{work}

      $P(x) = $ \ans{$23x - 80$} \qquad \textbf{(b)} $P(6) = $ \ans{$58$}

      \textbf{(c)} State the domain of $P$ in this context, and say why it is not all real numbers.
      \writespace{1.1cm}{Whole numbers of lawns from 0 to 12: $0 \le x \le 12$. You cannot mow half a lawn or a negative one, and the service can handle at most 12 in a day.}

\end{enumerate}

\vspace{0.1in}

\begin{headlinebox}{lilac}
\textbf{\color{plum}Part D --- Find the error, and the headline}
\end{headlinebox}

\begin{enumerate}[leftmargin=1.9em, itemsep=8pt, start=8]

\item A classmate writes: ``$f(x) = x^2 - 1$ and $g(x) = x - 1$, so
      $\left(\dfrac{f}{g}\right)(x) = x + 1$. That is a line, and lines are defined everywhere,
      so the domain is all real numbers.'' Find the error and give the correct domain.
      \writespace{1.5cm}{The quotient is undefined at $x = 1$, since $g(1) = 0$. The domain is every real number except $1$. The simplified line $x+1$ is only equal to $f/g$ where $f/g$ was defined in the first place.}

\item In your own words, explain how to find the domain of a combination of two functions. Use
      the words \emph{overlap} and \emph{zero}, and say why the simplified formula is never the
      place to read it from.
      \writespace{1.5cm}{Start from the two original functions, not the simplified formula: take the overlap of their domains, and if you divided, throw out every input where the bottom function is zero. Simplifying can hide a zero that was already excluded.}

\end{enumerate}

\vspace{0.08in}

\boxguard[8]
\begin{spiralbox}
\small
\textbf{Next lesson (1.3, Function Composition and Decomposition):} today two machines were fed
the \emph{same} input and their outputs were combined afterward. Next we chain them instead ---
the output of one function becomes the \textbf{input} of the next, written $f(g(x))$. Look back
at problem 6: if the shop's shipping charge depended on the total cost $T$, you would feed
$T(x)$ straight into the shipping rule rather than adding the two.
\end{spiralbox}

\end{document}
